%% Beginning of file 'sample631.tex'
%%
%% Modified 2022 May
%%
%% This is a sample manuscript marked up using the
%% AASTeX v6.31 LaTeX 2e macros.
%%
%% AASTeX is now based on Alexey Vikhlinin's emulateapj.cls
%% (Copyright 2000-2015).  See the classfile for details.

%% using aastex version 6.3
\documentclass[twocolumn]{aastex631}

\newcommand{\vdag}{(v)^\dagger}
\newcommand\aastex{AAS\TeX}
\newcommand\latex{La\TeX}

\usepackage{comment}
\usepackage{adjustbox}
\usepackage{amsmath}
\usepackage{breakurl}
\usepackage{listings}
\usepackage{xcolor}
\usepackage{placeins}

\def\UrlBreaks{\do\/\do-}

\definecolor{codegreen}{rgb}{0,0.6,0}
\definecolor{codegray}{rgb}{0.5,0.5,0.5}
\definecolor{codepurple}{rgb}{0.58,0,0.82}
\definecolor{backcolour}{rgb}{0.95,0.95,0.92}

\lstdefinestyle{mystyle}{
    backgroundcolor=\color{backcolour},
    commentstyle=\color{codegreen},
    keywordstyle=\color{blue},
    numberstyle=\tiny\color{codegray},
    stringstyle=\color{codepurple},
    basicstyle=\ttfamily\footnotesize,
    breakatwhitespace=false,
    breaklines=true,
    captionpos=b,
    keepspaces=true,
    numbers=left,
    numbersep=5pt,
    showspaces=false,
    showstringspaces=false,
    showtabs=false,
    tabsize=2
}
\lstset{style=mystyle}

\pagenumbering{arabic}

\begin{document}

\title{ASTR 5140: Project 3}

\author[0009-0004-8200-9910]{Winston Zhang}
\affiliation{University of Virginia, Department of Astronomy \\
530 McCormick Rd \\
Charlottesville, VA 22904, USA}

\author{Nitya Kallivayalil}
\affiliation{University of Virginia, Department of Astronomy \\
530 McCormick Rd \\
Charlottesville, VA 22904, USA}

\clearpage
\pagenumbering{arabic}

\section{Black Hole Scaling Relations in IllustrisTNG}

\subsection{Step 1: Halo Mass vs. Black Hole Mass}

\textcolor{blue}{\textbf{Prompt:} After identifying the data columns you wish to use, plot the relationship between dark matter halo mass (x-axis) and black hole mass (y-axis) for both subhalo (SubFind) and group (FoF) quantities.
These plots must be included in the assignment you turn in, in addition to a short description of which
data columns you used and why. If you performed any additional data cleaning or alteration you should
summarize it here as well. Given the data volume you should think about how to present this data most
clearly.}

\textbf{Columns and Data Preparation:}
Masses were converted from simulation units to physical $M_\odot$ using the $10^{10}/h$ scaling given in the IllustrisTNG documentation. To enable log--log visualization, objects with non-finite or non-positive halo mass or black hole mass were excluded. Because the catalogs contain millions of entries, relations were visualized with a log-scaled \texttt{hexbin} density plot to avoid overplotting while retaining information about the number density in parameter space.

\textbf{SubFind (Subhalo) Relation:}

For the SubFind catalog, the total bound subhalo mass (\texttt{SubhaloMass}) was used as the halo mass proxy, while the total black hole mass within each subhalo (\texttt{SubhaloBHMass}) was used for the black hole mass. The SubFind catalog identifies gravitationally bound structures corresponding approximately to individual galaxies embedded within larger halos. This relation therefore probes the connection between black hole mass and galaxy-scale bound mass.

\begin{figure}
    \centering
    \includegraphics[width=\linewidth]{science images/step1-subfind.png}
    \caption{Black hole mass versus bound subhalo mass from the SubFind catalog. Each hexagon represents the logarithmic number density of objects in that region of parameter space.}
    \label{fig:step1-subfind}
\end{figure}

\textbf{FoF (Group) Relations:}

For the FoF catalog, the total black hole mass associated with each group (\texttt{GroupBHMass}) was used for the y-axis. Two different halo mass definitions were examined:

\begin{itemize}
    \item \texttt{GroupMass}: the total mass of all particles linked by the friends-of-friends algorithm.
    \item \texttt{Group\_M\_Crit200}: the spherical overdensity mass $M_{200c}$, defined as the mass enclosed within a radius where the mean density equals 200 times the critical density of the universe.
\end{itemize}

The use of both mass definitions allows comparison between a halo mass defined by FoF linking and a physically motivated spherical overdensity mass commonly used in observational and theoretical studies.

\begin{figure}
    \centering
    \includegraphics[width=\linewidth]{science images/step1-fof-groupmass.png}
    \caption{Black hole mass versus FoF group mass (\texttt{GroupMass}).}
    \label{fig:step1-fof-groupmass}
\end{figure}

\begin{figure}
    \centering
    \includegraphics[width=\linewidth]{science images/step1-fof-m200c.png}
    \caption{Black hole mass versus spherical overdensity halo mass $M_{200c}$ (\texttt{Group\_M\_Crit200}).}
    \label{fig:step1-fof-m200c}
\end{figure}

\subsection{Step 2: Physical Scaling Between Halo Mass and Black Hole Mass}
\textcolor{blue}{\textbf{Prompt:} Come up with a reasonable physical equation relating the black hole mass to the halo mass. You do not need to estimate physical constants, but you should explain why you chose the scaling that you
did. Hypothesize as to the physical mechanisms that could give rise to such a relationship between these
parameters.}

\textbf{Proposed Scaling Relation:}

\begin{equation}
M_{\rm BH} = A\, M_{\rm halo}^{\alpha}
\end{equation}

A power-law relation provides a natural description of the trends in Figures~\ref{fig:step1-subfind}--\ref{fig:step1-fof-m200c}: the relation is approximately linear in log--log space, consistent with $M_{\rm BH}\propto M_{\rm halo}^{\alpha}$. Similar scalings between black hole mass and host/halo properties are commonly reported in observational and theoretical work \citep{2002ApJ...578...90F}.

\subsection{Step 3: SubFind vs. FoF Comparison}
\textcolor{blue}{\textbf{Prompt:} Determine if there is a difference in the relation between the SubFind and FoF catalogs. Form a hypothesis as to why there should or should not be a difference and include it in your write-up.}

At low halo masses, the SubFind relation exhibits noticeably larger scatter than the FoF relation. This is likely a consequence of the black hole seeding prescription in the simulation: black holes are introduced only once halos exceed a threshold mass and begin with a fixed seed mass. Many low-mass subhalos therefore host either no black hole or a newly seeded one, producing substantial dispersion.

Additional scatter arises from environmental effects such as tidal stripping. Satellite subhalos can lose dark matter mass while retaining their central black holes, further broadening the relation at the low-mass end.

In contrast, FoF halos combine the mass of multiple galaxies and their black holes, which smooths the overall trend and reduces stochasticity, particularly at higher masses. At the massive end, both relations become tighter as black hole growth is dominated by mergers and sustained gas accretion, leading to a more regulated scaling with halo mass.

Overall, the SubFind relation primarily traces black hole--galaxy co-evolution, while the FoF relation reflects the connection between black holes and the assembly history of their host dark matter halos.

\subsection{Step 4: Stellar Mass vs. Black Hole Mass}
\textcolor{blue}{\textbf{Prompt:} Plot the relation between galaxy stellar mass and black hole mass and comment on it briefly. Is it tighter or more scattered than the halo mass -- black hole mass relation? Is the correlation what you
expected? And so on. You may provide other supporting plots here if you wish.}}

A plot was created using the stellar mass of each galaxy, taken from the SubhaloMassType column (masked to type 4), and SubhaloBHMass for the black hole mass. This is shown in Figure~\ref{fig:step4-stellar-vs-bh}.

\begin{figure}
\centering
\includegraphics[width=\linewidth]{science images/step4-stellar-vs-bh.png}
\caption{Black hole mass vs. stellar mass for SubFind galaxies.}
\label{fig:step4-stellar-vs-bh}
\end{figure}

The stellar mass--black hole mass relation shows a clear positive correlation and is generally as tight as the halo mass--black hole mass relation. This is expected because stellar mass directly traces the baryonic component of galaxies and their star formation histories, both of which are closely linked to black hole growth through gas accretion and feedback processes.

At low stellar masses the relation broadens and exhibits a visible low-$M_{\rm BH}$ floor, consistent with the simulation's black hole seeding prescription.

\subsection{Step 5: Additional Correlated Quantity}
\textcolor{blue}{\textbf{Prompt:} Explore the data columns to identify another column that is correlated with black hole mass. Plot the relation and hypothesize as to the origin of the correlation.}}

Several parameters were plotted against black hole mass: velocity dispersion, star formation rate, gas mass, and maximum circular velocity. These are grouped in Figure~\ref{fig:step5-param-plots}.

\begin{figure*}
\centering
\includegraphics[width=\linewidth]{science images/step5-param-plots.png}
\caption{Black hole mass vs. selected correlated properties. Top left to bottom right: velocity dispersion, star formation rate, gas mass, and maximum circular velocity.}
\label{fig:step5-param-plots}
\end{figure*}

Velocity dispersion and maximum circular velocity show the strongest correlations with black hole mass. Both quantities trace the depth of the gravitational potential in the central regions of galaxies. A deeper potential well promotes gas inflow toward the galaxy center, enabling sustained black hole growth through accretion and mergers. This leads naturally to a scaling between black hole mass and dynamical measures of the host system.

In contrast, the correlations with star formation rate and gas mass are weaker. Star formation rate is a time-dependent quantity that reflects current activity rather than the integrated evolutionary history of the galaxy, and therefore does not track black hole growth as tightly. Similarly, total gas mass does not necessarily indicate how efficiently gas reaches the central black hole, as gas can remain in extended disks or be removed through feedback processes. Additionally, feedback from active galactic nuclei can suppress star formation in massive systems, producing quenched galaxies with large black holes but low star formation rates. This further weakens any direct correlation between star formation and black hole mass.

\subsection{Step 6: Limitations and Caveats}
\textcolor{blue}{\textbf{Prompt:} Comment on any limitations of the data products that you believe may impact the above
relations. If your work were to be published in a journal paper, what concerns would you have?}}

Several limitations of the IllustrisTNG data products may impact the scaling relations presented above.

First, the black hole seeding prescription introduces artificial floors in many of the relations. Black holes are inserted with a fixed seed mass once halos exceed a threshold mass, producing a population of galaxies with similar black hole masses at the low-mass end. This results in horizontal features in log--log space and increased scatter that are numerical artifacts rather than physical predictions. Any quantitative interpretation of the low-mass regime must therefore be treated with caution.

Second, the finite resolution of the simulation limits the reliability of results for low-mass galaxies. Small halos may be only marginally resolved, leading to uncertain stellar masses, velocity dispersions, and gas properties. Resolution effects can artificially broaden correlations or suppress intrinsic trends.

Third, filtering for $M_{\rm BH} > 0$ introduces a selection bias by excluding halos that do not host seeded black holes. This biases the sample toward systems above the seeding threshold and affects occupation fractions at low mass. As a result, derived scaling relations do not represent the full halo population but only black hole--hosting systems.

Finally, structural definitions differ between FoF groups and SubFind subhalos. FoF masses are based on a percolation algorithm, while $M_{200c}$ is defined via a spherical overdensity criterion. These definitions can diverge during mergers or in dense environments, leading to small but non-negligible differences in halo mass estimates. Such systematic uncertainties could affect slope measurements or normalization if one were to perform precise quantitative fits.

If this work were to be published, additional robustness tests would be necessary. These would include resolution convergence checks, sensitivity tests to the seeding prescription, and comparisons across different halo mass definitions.

% ============================================================
\section{Graduate Section (ASTR 5140 Only)}
% ============================================================

\subsection{1. Black Hole Mass--Velocity Dispersion Relation}
\textcolor{blue}{\textbf{Prompt:} Make cuts on the halo catalog to identify a list of subhalos with well-defined black hole masses and velocity dispersions (SubFind field [`SubhaloVelDisp`]). Summarize your cuts and plot this relation. Fit a model of your own design to this relation via least-squares or other similar technique. Comment on your result.}}

\textbf{Data Selection and Cuts:}

Subhalos were selected from the SubFind catalog using their velocity dispersions
(\texttt{SubhaloVelDisp}) and black hole masses (\texttt{SubhaloBHMass}).
Both quantities were converted to physical units and filtered to ensure reliable measurements.
As with all previous plots, objects were required to have finite, positive values of velocity dispersion and black hole mass.

Additional cuts were applied to isolate well-defined systems. Subhalos with
$M_{\rm BH} < 10^{6}\,M_\odot$ were excluded to avoid the artificial mass floor imposed by the
black hole seeding prescription in the simulation. A lower limit of
$\sigma > 20\ \mathrm{km\ s^{-1}}$ was also applied to remove poorly resolved, low-mass systems.
After these cuts, \(\sim 3.7\times10^{4}\) galaxies remained for analysis.

\textbf{Relation Between Black Hole Mass and Velocity Dispersion:}

Figure~\ref{fig:grad-step1-mbh-sigma} shows the $M_{\rm BH}$--$\sigma$ relation for the selected subhalos, with the best-fit power law overplotted.
The distribution exhibits a clear positive correlation, with more massive black holes residing in
galaxies with larger velocity dispersions. This trend is expected because velocity dispersion traces
the depth of the central gravitational potential, which regulates both black hole growth and bulge
formation.

\begin{figure}
\centering
\includegraphics[width=\linewidth]{science images/grad-step1-mbh-vs-sigma.png}
\caption{Black hole mass as a function of velocity dispersion for selected SubFind subhalos.
The red line shows the best-fit power-law relation obtained via least-squares fitting in log space.}
\label{fig:grad-step1-mbh-sigma}
\end{figure}

\textbf{Model Fit:}

The relation was modeled using a power-law scaling of the form
\begin{equation}
\log_{10}(M_{\rm BH}) = \alpha + \beta \log_{10}(\sigma/200),
\end{equation}
where $\sigma$ is normalized to $200\ \mathrm{km\ s^{-1}}$.

A least-squares fit in log space yielded
\begin{align}
\beta &\approx 4.2, \\
A &\approx 1.6 \times 10^{9}\,M_\odot,
\end{align}
with a strong positive correlation (Pearson $r$ close to unity). The fitted relation is overplotted
in Figure~\ref{fig:grad-step1-mbh-sigma}.

\textbf{Interpretation:}

The fitted slope supports the
interpretation that black hole growth is closely coupled to the dynamical state of the host galaxy
bulge. Feedback from accretion can regulate gas inflows and star formation, producing a
self-regulated scaling between black hole mass and central velocity dispersion.

Scatter in the relation arises from several factors, including differences in merger histories,
gas accretion rates, and the stochastic nature of feedback processes. Residual effects of the
seeding prescription and finite numerical resolution may also contribute, particularly at the
low-dispersion end.

Overall, the simulation reproduces a realistic $M_{\rm BH}$--$\sigma$ scaling, indicating that the
subgrid prescriptions for black hole growth and feedback capture the essential physics governing
black hole--galaxy co-evolution.

\subsection{2. AGN Luminosities}
\textcolor{blue}{\textbf{Prompt:} Assume all black holes accrete at their Eddington luminosities and compute the corresponding AGN luminosities for the galaxies you selected from IllustrisTNG.}}

\subsubsection{2a. Average Luminosity vs. Halo Mass}

\textcolor{blue}{\textbf{Prompt:} Calculate and plot the average AGN luminosity as a function of halo mass.}}

For each selected SubFind subhalo, the host FoF halo was identified using the group index
\texttt{SubhaloGrNr}, which maps each subhalo to an entry in the FoF catalog. Host halo masses were
taken from \texttt{Group\_M\_Crit200} (converted to physical units), since \(M_{200c}\) is a standard
spherical-overdensity halo mass definition used in galaxy formation studies.

Assuming Eddington-limited accretion, the AGN luminosity was computed from the black hole mass as
\begin{equation}
L_{\rm Edd} = 1.26 \times 10^{38}\left(\frac{M_{\rm BH}}{M_\odot}\right)\ \mathrm{erg\ s^{-1}},
\end{equation}
using \texttt{SubhaloBHMass} for \(M_{\rm BH}\). To estimate the average trend,
\(\langle L_{\rm Edd}\rangle\) was computed in logarithmically spaced bins of \(M_{200c}\).

Figure~\ref{fig:grad-step2a-ledd-m200c} shows \(L_{\rm Edd}\) as a function of host halo mass, with a
binned mean trend overlaid. As expected, the mean Eddington luminosity increases with halo mass,
reflecting the underlying correlation between black hole mass and the depth of the host potential well.

At the highest halo masses, however, the mean trend begins to flatten. This behavior arises because
halo mass can continue to grow through dark matter accretion and mergers even after black hole growth
becomes self-regulated. In massive systems, AGN feedback suppresses gas inflow and star formation,
limiting further black hole accretion. As a result, \(M_{\rm BH}\) (and therefore \(L_{\rm Edd}\))
increases more slowly than \(M_{200c}\), weakening the scaling at the massive end. The flattening
therefore reflects the decoupling of halo assembly from central black hole growth in the most massive
systems.

\begin{figure}[ht]
\centering
\includegraphics[width=\linewidth]{science images/grad-step2a-Ledd-vs-m200c.png}
\caption{Eddington luminosity \(L_{\rm Edd}\) as a function of host halo mass \(M_{200c}\) for the
selected subhalos. The solid line shows the binned mean \(\langle L_{\rm Edd}\rangle\) in
log-spaced halo-mass bins.}
\label{fig:grad-step2a-ledd-m200c}
\end{figure}

\subsubsection{2b. Luminosity Distribution}
\textcolor{blue}{\textbf{Prompt:} Plot the distribution (histogram) of AGN luminosities.}}

Using the same luminosities as before, a histogram was created to show the distribution of AGN luminosities, shown in Figure~\ref{fig:grad-step2b-luminosity}. As expected, this is a projection of Figure~\ref{fig:grad-step2a-ledd-m200c} along the vertical axis.

\begin{figure}
\centering
\includegraphics[width=\linewidth]{science images/grad-step2b-luminosity.png}
\caption{Distribution of Eddington AGN luminosities.}
\label{fig:grad-step2b-luminosity}
\end{figure}

\subsubsection{2c. Fraction of Halos Above Threshold}
\textcolor{blue}{\textbf{Prompt:} What fraction of halos in TNG100-1 have Eddington AGN luminosities greater than $1.28\times10^{46}\ \mathrm{erg\ s^{-1}}$?}}

Of the 36,883 selected subhalos, there were 3,416 above the threshold of $1.28\times10^{46}\ \mathrm{erg\ s^{-1}}$. This corresponds to a fraction of 0.0926 (9.26\%).

\subsection{3. Using TNG Accretion Rates}

\subsubsection{3a. Accretion Rate Comparison}
\textcolor{blue}{\textbf{Prompt:} Make a plot comparing the black hole accretion rates in the TNG catalogs to the Eddington mass accretion rate as a function of halo mass. Discuss whether the TNG results are consistent with Eddington accretion.}}

Figure~\ref{fig:grad-step3a-accretion} compares the catalog black hole accretion rates,
\(\dot{M}_{\rm TNG}\), to the corresponding Eddington mass accretion rates,
\(\dot{M}_{\rm Edd}\), as a function of host halo mass \(M_{200c}\).
The right panel shows the ratio \(\dot{M}_{\rm TNG}/\dot{M}_{\rm Edd}\),
with the dashed line marking exact Eddington accretion.

The median ratio is
\(\dot{M}_{\rm TNG}/\dot{M}_{\rm Edd} \sim 7 \times 10^{-4}\),
indicating that typical black holes in TNG accrete at only
$\sim 0.1\%$ of the Eddington rate.
Only a small fraction of objects approach or exceed the Eddington limit.

These results demonstrate that Eddington accretion represents an upper bound
rather than a typical state in the simulation.
Black hole growth in TNG is strongly self-regulated by
AGN feedback.
As a result, most black holes spend the majority of their time in
the sub-Eddington regime.

Overall, the TNG results are not consistent with the assumption that
black holes generically accrete at the Eddington rate.

\begin{figure*}
\centering
\includegraphics[width=\linewidth]{science images/grad-step3a-accretion-rate.png}
\caption{TNG vs. Eddington accretion rates.}
\label{fig:grad-step3a-accretion}
\end{figure*}

\subsubsection{3b. Luminosities from TNG Accretion}
\textcolor{blue}{\textbf{Prompt:} Let the efficiency of converting mass accretion into luminosity ($\epsilon r$ in the reading) be 0.1. Compute AGN luminosities for the galaxies you selected from IllustrisTNG using the black hole accretion rates in the TNG catalogs. Plot the average AGN luminosity calculated this way as a function of halo mass and compare to the same plot you made in 2a) above assuming Eddington mass accretion. Plot the distribution of AGN luminosities and compare to the plot from 2b) above. What fraction of halos in TNG100-1 have AGN luminosities greater than 1.28e46 ergs / s when using the TNG mass accretion rates? How does this compare to your answer from 2c) where you assumed Eddington accretion? Comment on how the assumption of Eddington accretion rates may affect predictions for the frequency of observing bright quasars.}}

Using the TNG black hole accretion rates (\texttt{SubhaloBHMdot}), AGN luminosities were computed assuming a radiative efficiency $\epsilon_r = 0.1$,
\begin{equation}
L_{\rm TNG} = \epsilon_r \dot{M}_{\rm BH} c^2.
\end{equation}
Accretion rates were converted to cgs units to evaluate the luminosity in $\mathrm{erg\ s^{-1}}$. Host halo masses were again taken from $M_{200c}$.

Figure~\ref{fig:grad-step3b-tng-accretion} shows $L_{\rm TNG}$ as a function of halo mass, with the binned mean $\langle L_{\rm TNG}\rangle$ overlaid. For comparison, the mean Eddington luminosity from Step~2a is also shown.

The difference between the two curves is substantial. The median luminosity computed from the TNG accretion rates is
\[
\mathrm{median}(L_{\rm TNG}) \approx 3.4 \times 10^{41}\ \mathrm{erg\ s^{-1}},
\]
whereas the median Eddington luminosity is
\[
\mathrm{median}(L_{\rm Edd}) \approx 3.1 \times 10^{44}\ \mathrm{erg\ s^{-1}}.
\]
The median ratio is
\[
\frac{L_{\rm TNG}}{L_{\rm Edd}} \sim 7 \times 10^{-4},
\]
indicating again that typical black holes accrete at 0.1\% of their Eddington luminosity.

The luminosity distribution derived from TNG accretion rates is therefore shifted downward by three orders of magnitude relative to the Eddington-based distribution from Step~2b. High-luminosity systems are correspondingly rare.

This comparison demonstrates that assuming Eddington-limited accretion significantly overpredicts the frequency of bright quasar-like systems. In the simulation, most black holes spend the majority of their time in low-accretion states regulated by feedback and limited gas supply. Eddington accretion therefore represents an upper envelope rather than a typical operating state. Consequently, models that assume universal Eddington accretion will substantially overestimate the abundance of luminous AGN.

\begin{figure}
\centering
\includegraphics[width=\linewidth]{science images/grad-step3b-tng-accretion-lum.png}
\caption{AGN luminosity from TNG accretion vs halo mass.}
\label{fig:grad-step3b-tng-accretion}
\end{figure}

\subsubsection{3c. Luminosity Distribution}
\textcolor{blue}{\textbf{Prompt:} Plot the distribution of AGN luminosities and compare to the plot from 2b.}}

The distributions of AGN luminosities were computed for the same set of selected subhalos used in the previous steps. Two luminosity estimates were compared:

\begin{itemize}
    \item $L_{\rm Edd}$: luminosities assuming Eddington-limited accretion, computed from black hole mass.
    \item $L_{\rm TNG}$: luminosities derived from the TNG black hole accretion rates using
    $L_{\rm TNG} = \epsilon_r \dot{M}_{\rm BH} c^2$
    with radiative efficiency $\epsilon_r = 0.1$.
\end{itemize}

Figure~\ref{fig:grad-step3c-tng-vs-edd-accretion} shows the probability density distributions of $\log_{10} L$ for both cases using logarithmic bins. The vertical dashed line marks the quasar threshold of $1.28\times10^{46}\ \mathrm{erg\ s^{-1}}$.

The two distributions differ substantially. The Eddington-based luminosities are systematically higher, peaking near $\log_{10}(L/\mathrm{erg\ s^{-1}})\sim 44$--46, while the TNG-derived luminosities are shifted to much lower values, with a median several orders of magnitude smaller. The median ratio between the two is about $7\times10^{-4}$, indicating that most black holes in the simulation accrete far below the Eddington limit. Physically, this reflects the fact that sustained Eddington accretion is rare.

\begin{figure}
\centering
\includegraphics[width=\linewidth]{science images/grad-step3c-tng-vs-edd-accretion.png}
\caption{Distribution of AGN luminosities from TNG accretion.}
\label{fig:grad-step3c-tng-vs-edd-accretion}
\end{figure}

\subsubsection{3d. Bright Quasar Fraction and Implications}

\textcolor{blue}{\textbf{Prompt:} What fraction of halos in TNG100-1 have AGN luminosities greater than 1.28e46 ergs/s
when using the TNG mass accretion rates? How does this compare to your answer from
2c) where you assumed Eddington accretion? Comment on how the assumption of
Eddington accretion rates may affect predictions for the frequency of observing bright
quasars.}}

Using the luminosity-valid selected sample ($N = 36{,}883$ subhalos), we computed the fraction of
objects exceeding a fixed quasar-like luminosity threshold of $1.28\times10^{46}\ \mathrm{erg\ s^{-1}}$. When assuming Eddington-limited accretion (Step 2c), $3{,}416$ objects exceeded the threshold,
corresponding to a fraction
\begin{equation}
f_{\rm Edd} = \frac{3416}{36883} \approx 0.0926 \quad (9.26\%).
\end{equation}
In contrast, when luminosities were computed from the \textit{actual} TNG accretion rates with
$\epsilon_r=0.1$ (Step 3b),
\begin{equation}
L_{\rm TNG} = \epsilon_r \dot{M}_{\rm BH} c^2,
\end{equation}
no objects in the same sample exceeded $L_{\rm thresh}$, giving
\begin{equation}
f_{\rm TNG} \approx 0.
\end{equation}

This comparison shows that assuming Eddington accretion dramatically overpredicts the incidence of
bright quasars. In the simulation, most black holes accrete at strongly sub-Eddington rates, so the high-luminosity
tail is greatly suppressed relative to an Eddington-based estimate. Physically, this reflects that
Eddington accretion is not sustained for most systems: gas supply limits, feedback self-regulation,
and duty-cycle effects make near-Eddington phases rare and short-lived. Therefore, using Eddington
accretion as a default assumption can bias predictions toward too many observable, extremely luminous
AGN, whereas using the catalog accretion rates produces a much lower expected frequency of bright
quasar events at a given snapshot in time.


\bibliography{sources}{}
\bibliographystyle{aasjournal}

\end{document}